\documentclass[11pt,letterpaper]{article}
\usepackage[margin=0.76in,top=0.68in,bottom=0.62in]{geometry}
\usepackage[T1]{fontenc}
\usepackage[utf8]{inputenc}
\usepackage{lmodern}
\usepackage[portuguese]{babel}
\usepackage[table]{xcolor}
\usepackage{array,tabularx}
\usepackage{enumitem}
\usepackage{fancyhdr}
\usepackage[most]{tcolorbox}
\usepackage{microtype}
\usepackage[hidelinks]{hyperref}

\definecolor{ideaorange}{HTML}{C45400}
\definecolor{browntext}{HTML}{713B19}
\definecolor{darkgray}{HTML}{373737}
\definecolor{softgray}{HTML}{F1F3F5}
\definecolor{cream}{HTML}{FFF2E6}
\definecolor{personacell}{HTML}{F5E5D7}
\definecolor{headergray}{HTML}{777777}

\pagestyle{fancy}
\fancyhf{}
\fancyhead[R]{\small\color{headergray}RepHouse\quad|\quad Ideação do Projeto}
\fancyfoot[C]{\small\color{headergray}Página \textcolor{black}{\thepage}}
\renewcommand{\headrulewidth}{0pt}
\renewcommand{\footrulewidth}{0pt}
\setlength{\headheight}{14pt}
\setlength{\parindent}{0pt}
\setlength{\parskip}{5pt}
\raggedbottom

\newcommand{\ideasection}[1]{%
  \vspace{4pt}{\color{ideaorange}\sffamily\bfseries\fontsize{17}{20}\selectfont #1\par}\vspace{7pt}}
\newcommand{\ideasub}[1]{%
  \vspace{5pt}{\color{browntext}\sffamily\bfseries\fontsize{12}{14}\selectfont #1\par}\vspace{2pt}}
\newcommand{\minisub}[1]{%
  \vspace{4pt}{\color{darkgray}\sffamily\bfseries\fontsize{10.5}{12}\selectfont #1\par}\vspace{1pt}}
\newcommand{\tightlist}{\begin{itemize}[leftmargin=1.55em,itemsep=0pt,topsep=1pt,parsep=0pt]}
\newcommand{\tightlistend}{\end{itemize}}
\newcolumntype{Y}{>{\raggedright\arraybackslash}X}

\tcbset{
  ideabox/.style={sharp corners,boxrule=0pt,left=8pt,right=8pt,top=6pt,bottom=6pt,
    before skip=5pt,after skip=8pt},
  graybox/.style={ideabox,colback=softgray,colframe=softgray},
  creambox/.style={ideabox,colback=cream,colframe=cream}
}

\begin{document}
\fontsize{10.8}{13}\selectfont

% Página 1
\vspace*{1.45in}
\begin{center}
  {\color{ideaorange}\bfseries\fontsize{31}{35}\selectfont RepHouse\par}
  \vspace{20pt}
  {\color{darkgray}\bfseries\fontsize{26}{30}\selectfont Ideação do Sistema\par}
  \vspace{12pt}
  {\color{headergray}\fontsize{15}{19}\selectfont Planejamento do projeto do zero - Itens 1 a 12\par}
\end{center}
\vspace{55pt}
\begin{tcolorbox}[creambox]
  \textcolor{ideaorange}{\textbf{Objetivo:}} Estruturar a ideia do RepHouse antes do desenvolvimento, partindo do problema,
  público-alvo, dores, brainstorming, proposta de solução, MVP e funcionalidades futuras.
\end{tcolorbox}
\vspace{32pt}
\begin{center}
  {\color{headergray}\itshape Documento de ideação e definição inicial do produto}
\end{center}

\newpage
% Página 2
\ideasection{1. Introdução}
O RepHouse é uma proposta de sistema web voltado para a organização e administração de
repúblicas estudantis e residências compartilhadas.

A ideia surgiu a partir da percepção de que pessoas que dividem uma residência frequentemente
precisam controlar despesas, tarefas domésticas, moradores, responsabilidades e informações da
casa utilizando diferentes ferramentas, como grupos de mensagens, planilhas, anotações e
aplicativos separados.

Essa descentralização pode gerar dificuldades de organização, conflitos entre moradores,
esquecimentos e falta de transparência.

Diante desse cenário, surgiu a proposta de desenvolver uma plataforma que concentre as
principais necessidades de gestão de uma república em um único ambiente digital.

\ideasection{2. Problema}
Pessoas que moram em repúblicas ou residências compartilhadas precisam organizar diferentes
aspectos do dia a dia, entre eles:
\tightlist
  \item despesas
  \item divisão de contas
  \item tarefas domésticas
  \item moradores
  \item responsabilidades
  \item informações importantes da residência
  \item organização financeira
  \item comunicação
\tightlistend
Normalmente essas atividades são realizadas utilizando diferentes meios. Por exemplo:
\begin{tcolorbox}[graybox]
WhatsApp -> comunicação\\
Planilha -> despesas\\
Papel -> tarefas\\
Aplicativo bancário -> pagamentos\\
Calendário -> compromissos
\end{tcolorbox}
Isso pode provocar:
\tightlist
  \item informações espalhadas
  \item contas esquecidas
  \item dificuldade para saber quem pagou
  \item divisão incorreta de despesas
  \item tarefas domésticas esquecidas
  \item conflitos entre moradores
  \item dificuldade de acompanhar responsabilidades
  \item falta de histórico
  \item ausência de uma visão geral da situação da residência
\tightlistend

\newpage
% Página 3
\ideasection{3. Pergunta norteadora}
\begin{tcolorbox}[creambox]
\textcolor{ideaorange}{\textbf{Pergunta central:}} Como podemos criar uma plataforma simples e centralizada que facilite a
organização de uma república e reduza conflitos relacionados a despesas, tarefas e
responsabilidades entre os moradores?
\end{tcolorbox}

\ideasection{4. Público-alvo}
O sistema será destinado principalmente a:
\tightlist
  \item estudantes que moram em repúblicas
  \item jovens que dividem apartamentos
  \item moradores de residências compartilhadas
  \item administradores de repúblicas
  \item grupos de pessoas que compartilham despesas e responsabilidades domésticas
\tightlistend

\ideasection{5. Personas}
Para compreender melhor os possíveis usuários do sistema, foram criadas personas fictícias que
representam perfis recorrentes do público-alvo.

\ideasub{Persona 1 - Administrador da República}
\begin{tabularx}{\textwidth}{|>{\bfseries\columncolor{personacell}}p{0.48\textwidth}|Y|}
\hline
Nome & Lucas \\
\hline
Idade & 24 anos \\
\hline
Perfil & Estudante universitário responsável por organizar a república. \\
\hline
\end{tabularx}

\minisub{Necessidades}
\tightlist
  \item cadastrar moradores
  \item controlar despesas
  \item acompanhar pagamentos
  \item distribuir tarefas
  \item visualizar a situação geral da casa
  \item remover moradores que saíram
  \item manter as informações organizadas
\tightlistend

\minisub{Problemas}
Atualmente Lucas utiliza:
\begin{tcolorbox}[graybox]
WhatsApp + Excel + anotações
\end{tcolorbox}

\newpage
% Página 4
Isso gera informações espalhadas e dificulta o controle.

\ideasub{Persona 2 - Morador}
\begin{tabularx}{\textwidth}{|>{\bfseries\columncolor{personacell}}p{0.48\textwidth}|Y|}
\hline
Nome & Gabriel \\
\hline
Idade & 21 anos \\
\hline
Perfil & Estudante que divide uma residência com outras pessoas. \\
\hline
\end{tabularx}

\minisub{Necessidades}
\tightlist
  \item saber quanto precisa pagar
  \item consultar despesas
  \item verificar suas tarefas
  \item acompanhar prazos
  \item saber quem mora na residência
  \item evitar cobranças incorretas
\tightlistend

\minisub{Problemas}
Gabriel muitas vezes não sabe:
\begin{tcolorbox}[graybox]
Quanto deve?\\
Quando precisa pagar?\\
Quem já pagou?\\
Qual tarefa precisa realizar?
\end{tcolorbox}

\ideasection{6. Levantamento das dores}
Durante a ideação foram identificadas algumas dores principais que a solução deverá tratar.

\ideasub{Gestão financeira}
Moradores precisam controlar contas como:
\tightlist
  \item aluguel
  \item energia
  \item água
  \item internet
  \item gás
  \item compras coletivas
  \item manutenção
  \item produtos de limpeza
\tightlistend
O problema ocorre quando essas informações são controladas manualmente ou ficam espalhadas
em diferentes ferramentas.

\newpage
% Página 5
\ideasub{Divisão de despesas}
Nem todas as despesas necessariamente precisam ser divididas igualmente. Por exemplo:
\begin{tcolorbox}[graybox]
Conta de energia: R\$ 400\\
4 moradores\\
R\$ 100 para cada
\end{tcolorbox}

Mas também podem existir situações específicas:
\begin{tcolorbox}[graybox]
Compra realizada somente por 3 moradores
\end{tcolorbox}

Portanto, o sistema deverá permitir diferentes formas de divisão.

\ideasub{Organização de tarefas}
Atividades como:
\tightlist
  \item limpar cozinha
  \item limpar banheiro
  \item retirar lixo
  \item comprar produtos
  \item organizar área comum
\tightlistend
podem ser esquecidas ou realizadas repetidamente pelas mesmas pessoas.

\ideasub{Gestão de moradores}
Também é necessário manter informações como:
\tightlist
  \item quem mora atualmente
  \item quem é administrador
  \item quando uma pessoa entrou
  \item quando saiu
  \item qual residência ela pertence
\tightlistend

\ideasection{7. Brainstorming}
A partir dos problemas identificados, foi realizada uma geração livre de ideias. O objetivo desta
etapa é explorar possibilidades antes de decidir o que realmente será desenvolvido.
\tightlist
  \item Cadastro de usuários
  \item Criação de república
  \item Convite de moradores
  \item Gestão de moradores
  \item Cadastro de despesas
  \item Divisão automática de despesas
  \item Controle de pagamentos
  \item Tarefas domésticas
  \item Calendário
\tightlistend

\newpage
% Página 6
\tightlist
  \item Notificações
  \item Dashboard
  \item Relatórios financeiros
  \item Ranking de tarefas
  \item Chat interno
  \item Mural de avisos
  \item Lista de compras
  \item Reserva de espaços
  \item Controle de patrimônio
  \item Histórico da república
  \item Integração PIX
  \item Notificação por e-mail
  \item Aplicativo mobile
\tightlistend
Nesse momento da ideação, o objetivo não é implementar todas as funcionalidades, mas
identificar possibilidades e organizar uma visão ampla do produto.

\ideasection{8. Agrupamento das ideias}
Após o brainstorming, as funcionalidades foram organizadas por áreas para facilitar o
entendimento e a priorização.

\ideasub{Usuários}
\tightlist
  \item Cadastro
  \item Login
  \item Perfil
  \item Recuperação de senha
\tightlistend

\ideasub{República}
\tightlist
  \item Criar república
  \item Editar informações
  \item Convidar moradores
  \item Remover moradores
  \item Definir administrador
\tightlistend

\ideasub{Financeiro}
\tightlist
  \item Cadastrar despesas
  \item Dividir despesas
  \item Registrar pagamentos
  \item Consultar pendências
  \item Visualizar histórico
\tightlistend

\ideasub{Tarefas}
\tightlist
  \item Criar tarefa
  \item Definir responsável
  \item Definir prazo
\tightlistend

\newpage
% Página 7
\tightlist
  \item Marcar como concluída
  \item Visualizar tarefas pendentes
\tightlistend

\ideasub{Comunicação}
\tightlist
  \item Mural de avisos
  \item Notificações
  \item Chat
\tightlistend

\ideasub{Gestão}
\tightlist
  \item Dashboard
  \item Relatórios
  \item Indicadores
  \item Histórico
\tightlistend

\ideasection{9. Proposta de solução}
A solução proposta é desenvolver uma aplicação web chamada RepHouse, responsável por
centralizar a gestão de uma residência compartilhada.

A plataforma permitiria que cada república possuísse seu próprio ambiente, com informações
separadas das demais.
\begin{tcolorbox}[graybox]
\begin{verbatim}
RepHouse
|
+-- República A
|   +-- Moradores
|   +-- Despesas
|   +-- Tarefas
|
+-- República B
|   +-- Moradores
|   +-- Despesas
|   +-- Tarefas
|
+-- República C
    +-- Moradores
    +-- Despesas
    +-- Tarefas
\end{verbatim}
\end{tcolorbox}
Os dados de uma república deverão permanecer separados dos dados das demais, garantindo
organização e privacidade.

\ideasection{10. Proposta de valor}
\begin{tcolorbox}[creambox]
\textcolor{ideaorange}{\textbf{Proposta de valor:}} Centralizar a administração de uma república em uma única plataforma
simples, organizada e acessível.
\end{tcolorbox}

\newpage
% Página 8
Em vez de utilizar várias ferramentas independentes:
\begin{tcolorbox}[graybox]
WhatsApp + Excel + papel + aplicativos diferentes
\end{tcolorbox}

o usuário poderá utilizar:
\begin{tcolorbox}[graybox]
RepHouse
\end{tcolorbox}

como ponto central da organização da residência.

\ideasection{11. MVP - Produto Mínimo Viável}
Não é recomendado desenvolver todas as ideias inicialmente. O primeiro lançamento deverá
conter somente as funcionalidades consideradas essenciais para validar a proposta.

\ideasub{MVP 1.0}
\minisub{Autenticação}
\tightlist
  \item criar conta
  \item fazer login
  \item logout
\tightlistend

\minisub{República}
\tightlist
  \item criar república
  \item editar informações
  \item visualizar república
\tightlistend

\minisub{Moradores}
\tightlist
  \item adicionar morador
  \item listar moradores
  \item remover morador
\tightlistend

\minisub{Despesas}
\tightlist
  \item cadastrar despesa
  \item informar valor
  \item informar vencimento
  \item dividir entre moradores
  \item visualizar despesas
  \item registrar pagamento
\tightlistend

\minisub{Tarefas}
\tightlist
  \item cadastrar tarefa
  \item selecionar responsável
  \item definir prazo
  \item concluir tarefa
\tightlistend

\newpage
% Página 9
\minisub{Dashboard}
\tightlist
  \item mostrar quantidade de moradores
  \item mostrar despesas pendentes
  \item mostrar valor total pendente
  \item mostrar próximos vencimentos
  \item mostrar tarefas pendentes
\tightlistend

\ideasection{12. Funcionalidades futuras}
Depois que o MVP estiver funcionando e validado com usuários, novas funcionalidades poderão
ser adicionadas de forma incremental.

\ideasub{Versão 2}
\tightlist
  \item Notificações
  \item Mural de avisos
  \item Lista de compras
  \item Relatórios financeiros
  \item Histórico de despesas
\tightlistend

\ideasub{Versão 3}
\tightlist
  \item Chat
  \item PIX
  \item Calendário
  \item Sistema de votação
  \item Reserva de áreas
\tightlistend

\ideasub{Versão 4}
\tightlist
  \item Aplicativo mobile
  \item Integrações bancárias
  \item Automação financeira
  \item Inteligência artificial
\tightlistend

\begin{tcolorbox}[creambox]
\textcolor{ideaorange}{\textbf{Encerramento desta etapa:}} Os itens 1 a 12 estabelecem a base da ideação do RepHouse. A partir
daqui, o projeto pode avançar para priorização, fluxos, prototipação, arquitetura e requisitos
detalhados.
\end{tcolorbox}

\end{document}
